% A0 portrait poster in the Laboratory house style: Beamer with beamerposter.
% Three columns of blocks, read top to bottom and left to right; a figure is
% drawn at the column width. Replace the pendulum with the work and keep the
% conventions. Build with build.py, which reads latexmkrc.
\documentclass{beamer}
\usepackage[size=a0, orientation=portrait, scale=1.4]{beamerposter}
\usetheme{LaboratoryPoster}

\title{Measuring the gravitational acceleration with a simple pendulum}
\author{Author Name}
\date{\today}
% The repository, notebook or data folder behind the poster; delete to omit.
\source{github.com/author/pendulum}

\begin{document}

\begin{frame}[t]
  \begin{columns}[t, totalwidth=\linewidth]
    \begin{column}{0.31\linewidth}
      \begin{block}{Question}
        A mass on a light string of length $L$ swings with a period that depends only on $L$ and the gravitational acceleration $g$ when the amplitude is small~\cite{nelson1986}. Timing the period at several lengths therefore measures $g$, and the spread of the measurements about the model tests the small-angle approximation.
      \end{block}

      \begin{block}{Model}
        The angle $\theta$ of the string from the vertical obeys
        \begin{equation}
          \frac{\dd^2\theta}{\dd t^2} = -\frac{g}{L} \sin\theta ,
          \label{eq:motion}
        \end{equation}
        which for small $\theta$ gives simple harmonic motion with period
        \begin{equation}
          T = 2\pi \sqrt{L/g} .
          \label{eq:period}
        \end{equation}
        At the \qty{3}{\degree} release used here, the finite amplitude lengthens the period by \num{1.7e-4}~\cite{nelson1986}.
      \end{block}

      \begin{block}{Method}
        A steel bob on a clamped string was released from rest. A light gate timed 20 oscillations at each of nine lengths, and each period is the mean of five runs with a standard error of \qty{1}{\milli\second}.
      \end{block}
    \end{column}

    \begin{column}{0.31\linewidth}
      \begin{block}{Results}
        Figure~\ref{fig:period} shows the measured periods with the fit of Equation~\eqref{eq:period}, and Table~\ref{tab:period} lists them at selected lengths.

        \begin{figure}
          \includegraphics[width=\linewidth]{figures/pendulum}
          \caption{Period against length with the fitted model (a), and residuals from the fit (b). Error bars are one standard error.}
          \label{fig:period}
        \end{figure}

        \begin{table}
          \caption{Measured and fitted periods at selected lengths.}
          \label{tab:period}
          \begin{tabular}{
            S[table-format=1.1]
            S[table-format=1.4(2)]
            S[table-format=1.4]
          }
            \toprule
            {\colhead{$L$}{\metre}} & {\colhead{$T$}{\second}} & {\colhead{$T_\mathrm{fit}$}{\second}} \\
            \midrule
            0.2 & 0.8983(10) & 0.8972 \\
            0.4 & 1.2692(10) & 1.2688 \\
            0.6 & 1.5548(10) & 1.5540 \\
            0.8 & 1.7937(10) & 1.7943 \\
            1.0 & 2.0057(10) & 2.0061 \\
            \bottomrule
          \end{tabular}
        \end{table}
      \end{block}
    \end{column}

    \begin{column}{0.31\linewidth}
      \begin{block}{Conclusion}
        A least-squares fit of $T = a \sqrt{L}$ gives
        \begin{equation*}
          g = \qty{9.809(4)}{\metre\per\second\squared} ,
        \end{equation*}
        within one standard error of the local value. The fit has $\chi^2 = 6.9$ for 8 degrees of freedom, so the timing error accounts for the scatter~\cite{taylor1997}, and the residuals show no trend with length.
      \end{block}

      \begin{block}{References}
        \bibliographystyle{apsrev4-2}
        \bibliography{refs}
      \end{block}
    \end{column}
  \end{columns}
\end{frame}

\end{document}
